\documentclass[11pt,a4paper]{article}
\usepackage[utf8]{inputenc}
\usepackage{amsmath,amssymb,amsthm}
\usepackage[margin=1in]{geometry}

\newtheorem{theorem}{Theorem}
\newtheorem{lemma}[theorem]{Lemma}
\newtheorem{definition}[theorem]{Definition}

\title{Problem 121: Disc Game Prize Fund}
\author{Project Euler Solution}
\date{}

\begin{document}
\maketitle

\section{Problem Statement}

A bag contains one red disc and one blue disc. Each turn, a red disc is added, then a disc is drawn at random and set aside. After $n = 15$ turns the player wins if more blue discs than red were drawn. Determine the maximum prize fund the banker should allocate.

\section{Mathematical Development}

\begin{definition}
At turn $k$ ($1 \le k \le n$), the bag contains 1 blue and $k$ red discs, for $k + 1$ total.
\end{definition}

\begin{theorem}[Turn probabilities]
$\Pr(\text{blue at turn } k) = \frac{1}{k+1}$ and $\Pr(\text{red at turn } k) = \frac{k}{k+1}$. The turns are independent.
\end{theorem}

\begin{proof}
The bag composition at turn $k$ is deterministic ($k + 1$ discs, exactly 1 blue). By a symmetry argument on the marginal draw probabilities, each turn's outcome is independent with the stated distribution.
\end{proof}

\begin{theorem}[Winning probability]
\[
P(\text{win}) = \frac{N}{(n+1)!}, \quad N = \sum_{\substack{T \subseteq \{1,\ldots,n\} \\ |T| \le \lfloor n/2 \rfloor}} \prod_{k \in T} k,
\]
where $T$ indexes turns on which red is drawn.
\end{theorem}

\begin{proof}
For an outcome where blue is drawn on turns in $S$ and red on turns in $T = \{1,\ldots,n\} \setminus S$:
\[
\Pr(S) = \prod_{k \in S} \frac{1}{k+1} \cdot \prod_{k \in T} \frac{k}{k+1} = \frac{\prod_{k \in T} k}{\prod_{k=1}^{n}(k+1)} = \frac{\prod_{k \in T} k}{(n+1)!}.
\]
The player wins when $|S| > n/2$, i.e., $|T| \le \lfloor n/2 \rfloor$. Summing numerators over qualifying $T$ yields $N$.
\end{proof}

\begin{lemma}[Elementary symmetric polynomial recurrence]
Let $e(i,j)$ be the sum of products of all $j$-element subsets of $\{1,\ldots,i\}$. Then
\[
e(i,j) = e(i-1,j) + i \cdot e(i-1,j-1), \quad e(0,0) = 1,
\]
and $N = \sum_{j=0}^{\lfloor n/2 \rfloor} e(n,j)$.
\end{lemma}

\begin{proof}
Partition $j$-element subsets of $\{1,\ldots,i\}$ by membership of $i$: those excluding $i$ contribute $e(i-1,j)$; those including $i$ contribute $i \cdot e(i-1,j-1)$. The sum $\sum_{j=0}^{\lfloor n/2\rfloor} e(n,j)$ collects all products $\prod_{k \in T} k$ over subsets $T$ of size at most $\lfloor n/2 \rfloor$.
\end{proof}

\begin{theorem}[Maximum prize fund]
The maximum prize fund is $\lfloor (n+1)!/N \rfloor$.
\end{theorem}

\begin{proof}
The player pays \$1. For prize $M$, expected payout is $M \cdot N/(n+1)!$. Non-negative expected profit requires $M \le (n+1)!/N$, so the maximum integer prize is $\lfloor (n+1)!/N \rfloor$.
\end{proof}

\section{Algorithm}

\begin{enumerate}
\item Initialize $dp[0] = 1$, $dp[j] = 0$ for $1 \le j \le \lfloor n/2 \rfloor$.
\item For $i = 1$ to $n$: for $j = \min(i, \lfloor n/2 \rfloor)$ down to $1$: update $dp[j] \mathrel{+}= i \cdot dp[j-1]$.
\item Compute $N = \sum_{j=0}^{\lfloor n/2 \rfloor} dp[j]$.
\item Return $\lfloor (n+1)!/N \rfloor$.
\end{enumerate}

\section{Complexity Analysis}

\begin{itemize}
\item \textbf{Time:} $O(n \cdot \lfloor n/2 \rfloor) = O(n^2)$, where $n = 15$.
\item \textbf{Space:} $O(n)$ for the rolling DP array.
\end{itemize}

\section{Answer}

\[
\boxed{2269}
\]

\end{document}
